\documentclass{article}
\usepackage{filecontents}

% ==============================================================================
% CREATE CONFIGURATION FILE
% ==============================================================================
\begin{filecontents*}{exampleconf.lua}
-- Custom configuration file for refsuite package
-- This file overrides the default configuration

local CustomConfig = {
    hyperref = {
        enable = true,
        loadtime_opts = {},
        shortref = {
            {csname = 'picref', reftext = 'Fig.'},
            {csname = 'tabref', reftext = 'Table'},
            {csname = 'secref', reftext = 'Section'}
        },
        hypersetup = {
            breaklinks = true,
            hyperindex = true,
            colorlinks = true,
            hidelinks = false,
            unicode = true,
            pdfauthor = {'John Doe'},
            pdfsubject = {'RefSuite Package Demonstration'},
            pdfkeywords = {'LaTeX, LuaTeX, bibliography, indexes'},
            pdftitle = {'RefSuite Demonstration'},
            bookmarksopen = true,
            linktocpage = true,
            plainpages = false,
            pdfpagelabels = true,
            -- Link colors
            urlcolor = 'blue',
            linkcolor = 'red',
            filecolor = 'magenta',
            citecolor = 'green'
        }
    },
    
    -- Enable bibliography support
    bibtex = {
        enable = true,
        loadtime_opts = {backend = 'biber', style = 'numeric'},
        path = 'build',  -- build directory
        style = 'numeric',
        
        -- Configuration for multiple bibliographies
        bibliographies = {
            {
                name = 'main',
                title = 'References',
                files = 'main.bib',
            },
            {
                name = 'git',
                title = 'Online Resources and Repositories',
                files = {'online.bib'},
            }
        }
    },
    
    -- Enable index support
    indexes = {
        enable = true,
        loadtime_opts = {'makeindex'},
        
        -- Index configuration
        list = {
            -- Notation index
            {
                name = 'notation',
                title = 'List of Notations',
                intoc = true,
                columns = 2,
                columnsep = '20pt',
            },
            -- Main subject index
            {
                name = nil,
                title = 'Subject Index',
                intoc = true,
                columns = 3,
                columnsep = '15pt',
            }
        }
    }
}

return CustomConfig
\end{filecontents*}

% ==============================================================================
% CREATE BIBLIOGRAPHY FILES
% ==============================================================================
\begin{filecontents*}{main.bib}
@book{knuth1984texbook,
  author = {Knuth, Donald E.},
  title = {The {\TeX}book},
  year = {1984},
  publisher = {Addison-Wesley},
  address = {Reading, Massachusetts},
}

@article{lamport1994latex,
  author = {Lamport, Leslie},
  title = {{\LaTeX}: A Document Preparation System},
  year = {1994},
  publisher = {Addison-Wesley},
  edition = {2nd},
}

@book{mittelbach2004latex,
  author = {Mittelbach, Frank and Goossens, Michel},
  title = {The {\LaTeX} Companion},
  year = {2004},
  publisher = {Addison-Wesley},
  edition = {2nd},
}
\end{filecontents*}

\begin{filecontents*}{online.bib}
@online{torvalds2003git,
  author = {Torvalds, Linus},
  title = {Git: Fast Version Control System},
  year = {2003},
  url = {https://git-scm.com},
  urldate = {2024-01-15},
}

@online{ctan2024latex,
  author = {{CTAN Team}},
  title = {Comprehensive {\TeX} Archive Network},
  year = {2024},
  url = {https://ctan.org},
  urldate = {2024-01-15},
}

@online{stackexchange2024tex,
  author = {{Stack Exchange Community}},
  title = {{\TeX} - {\LaTeX} Stack Exchange},
  year = {2024},
  url = {https://tex.stackexchange.com},
  urldate = {2024-01-15},
}
\end{filecontents*}

% ==============================================================================
% LOAD REFSUITE PACKAGE
% ==============================================================================
\usepackage[configuration=build/exampleconf.lua]{refsuite}

% ==============================================================================
% DOCUMENT CONTENT
% ==============================================================================
\begin{document}

\title{Example of Using the RefSuite Package}
\author{John Doe}
\date{\today}
\maketitle

\tableofcontents

\section{Introduction}

This document demonstrates the use of the \texttt{refsuite} package for managing references, bibliography, and indexes. All necessary files are created inline using the \texttt{filecontents} package.

\subsection{References to Document Elements}

The package automatically creates commands for short references. For example:

\begin{itemize}
    \item Reference to a figure: \picref{fig:sample} or \Picref{fig:sample}
    \item Reference to a table: \tabref{tab:sample} or \Tabref{tab:sample}
    \item Reference to a section: \secref{sec:conclusion} or \Secref{sec:conclusion}
\end{itemize}

\section{Main Content}

\subsection{Figures}

\begin{figure}[htbp]
    \centering
    \fbox{\parbox{0.5\textwidth}{\centering Example figure}}
    \caption{Example figure with caption}
    \label{fig:sample}
\end{figure}

As shown in \picref{fig:sample}, the package allows convenient referencing of graphic materials.

\subsection{Tables}

\begin{table}[htbp]
    \centering
    \begin{tabular}{|c|c|}
        \hline
        Column 1 & Column 2 \\
        \hline
        Data 1 & Data 2 \\
        Data 3 & Data 4 \\
        \hline
    \end{tabular}
    \caption{Example table}
    \label{tab:sample}
\end{table}

According to the data in \tabref{tab:sample}, conclusions can be drawn.

\subsection{Using Bibliography}

In scientific works, it is often necessary to cite sources \cite{knuth1984texbook}. 
For online resources, a separate bibliography is used \cite{torvalds2003git}.

\subsection{Working with Indexes}

You can add entries to different indexes:

\index[notation]{$\alpha$} Alpha notation
\index[notation]{$\beta$} Beta notation
\index{LaTeX} Typesetting system
\index{bibliography} List of references
\index{index} Alphabetical index

\section{Conclusion}
\label{sec:conclusion}

The \texttt{refsuite} package provides a convenient way to configure the reference, bibliography, and index system in LuaLaTeX. This standalone example shows how to use all features in a single file.

% ==============================================================================
% BIBLIOGRAPHY
% ==============================================================================
% Using the created commands to print bibliographies
% \printbibliographymain - for main literature
% \printbibliographygit - for online resources

\printbibliographymain
\printbibliographygit

% ==============================================================================
% INDEXES
% ==============================================================================
% Using the created commands to print indexes
% \printindexnotation - for list of notations
% \printindex - for subject index (default)

\printindexnotation
\printindex

\end{document}
